\documentclass[11pt,a4paper]{article}
\usepackage[T1]{fontenc}
\usepackage{lmodern,amsmath,amssymb,geometry}
\geometry{margin=28mm}
\title{EPPA--III: source-local repairs for Claims 3.15--3.16}
\author{Validation supplement; manuscript unchanged}
\date{9 October 2026}
\begin{document}
\maketitle

\paragraph{Scope and necessary hypothesis.}
The statement of Lemma 3.12 is false without excluding homogeneous
witnesses: \(K_2\dot\cup K_1\subseteq 3K_2\) and
\(K_3\dot\cup K_2\dot\cup K_1\subseteq4K_3\) are counterexamples.
We work under the necessary additional hypothesis that \(H\) is
\emph{non-homogeneous}. We use Claim 3.14, whose regularity/double-counting
argument is valid once Claim 3.13 is in place. The original source is
not modified.

\section*{A short replacement for Claim 3.13}
Let \(P\) be any invariant partition of the vertex set of a
vertex-transitive graph into blocks of size \(m\).
The setwise stabiliser of each block \(\Pi\) acts transitively
on \(\Pi\), so \(H[\Pi]\) is vertex-transitive.
In a finite regular graph of positive degree on \(m\) vertices,
an independent set has size at most \(m/2\), by counting its
incident edges. Thus a vertex-transitive graph having a clique
of size \(>m/2\) must be complete.
In Case 1, different components of \(G\) lie in different
blocks: otherwise the EPPA action on an independent transversal
would force all of \(G\) into one block.
As \(|G|=bm/2\) and its occupied-block clique sizes are unequal,
some block contains \(>m/2\) vertices of \(G\), forming a clique.
Therefore that block, and by transitivity every block, is a clique.

\section*{Replacement for the proof of Claim 3.15}

Write \(b=|P|\), \(m=|\Pi|\), and let \(S\subseteq P\) be the
\(s\) blocks meeting \(G\); put \(k_i=|G\cap\Pi_i|>0\).
Since \(|G|=bm/2\) and the sizes \(k_i\) are unequal,
we have \(s>b/2\).

Let \(Q\) be the graph on \(P\) joining two blocks whenever
\(H\) has an edge between them. Since \(H\) is not a disjoint
union of cliques and all blocks induce cliques, \(Q\) is
a positive-degree vertex-transitive graph.
Permutations of an independent transversal of the
\(s\) maximal cliques of \(G\) extend to automorphisms
of \(H\), inducing every permutation of the occupied
blocks \(S\). Hence \(Q[S]\) is either complete or
edgeless. It cannot be edgeless, because in a positive-degree
regular graph an independent set contains at most half
the vertices. Thus \(Q[S]\) is complete.
If \(Q\) were not complete, its complement would be
a positive-degree vertex-transitive graph having the
independent set \(S\) of size \(>b/2\). Therefore \(Q\)
itself is complete.

By Claim 3.14 there are no interblock edges between two
blocks disjoint from \(G\). As \(Q\) is complete,
there is at most one such block, say \(U\).
If \(U\) exists, extending permutations of the
independent transversal preserves \(U\) and is
transitive on the occupied blocks; by varying the
representatives, it is also transitive on \(V(G)\).
Thus every vertex of \(G\) has the same number
\(q>0\) of neighbours in \(U\), and the edge count
\(e(H[\Pi_i,U])=k_iq\) is independent of \(i\).
This forces all \(k_i\) to be equal, a contradiction.
Consequently every block is occupied: \(b=s\).

Fix \(x\in G\cap\Pi_i\). Extending permutations of an
independent transversal which fix \(x\) and permute the
other representatives shows that \(x\) has the same
number of neighbours in each \(\Pi_j\) with \(j\ne i\).
Since \(H\) is vertex-transitive, this number is a
common \(r>0\) for every vertex of \(H\).
For distinct blocks \(\Pi_i,\Pi_j\), Claim 3.14 says
all edges between them go from \(G\cap\Pi_i\) to
\(\Pi_j\setminus G\), or vice versa. Counting edges
between \(G\cap\Pi_i\) and \(\Pi_j\setminus G\)
from both sides therefore yields
\[
      k_i r=(m-k_j)r,\qquad\text{so}\qquad k_i+k_j=m.
\]
If \(s\ge3\), these equations force all the \(k_i\)'s
equal (in fact \(m/2\)), contrary to assumption.
Hence \(s=b=2\), as required. In particular, there
are precisely two imprimitivity blocks.

\section*{The next paragraph: the cross edges form a matching}

Let \(B_1,B_2\) be these two cliques of size
\(m=n\), and let \(G_i=G\cap B_i\) have unequal
sizes \(k_i\), with \(k_1+k_2=m\).
Let \(J\) be the graph consisting only of the edges
of \(H\) between \(B_1\) and \(B_2\).
It is \(r\)-regular for the positive \(r\) above
and is invariant under \(\operatorname{Aut}(H)\).
By Claim 3.14, it is the disjoint union of
\(r\)-regular bipartite graphs
\[
 J_1=J[G_1,B_2\setminus G_2]
 \quad\text{and}\quad
 J_2=J[B_1\setminus G_1,G_2]
\]
on \(2k_1\) and \(2k_2\) vertices, respectively.
Suppose \(r\ge2\). Some two points of \(G_1\)
have a common neighbour in \(J_1\), so are in
the same connected component of \(J\).
EPPA extends all permutations of the clique \(G_1\);
hence every two points of \(G_1\) lie in the same
component. As \(J_1\) has no isolated vertices,
\(J_1\) is connected. The same holds for \(J_2\).
But \(J\) is vertex-transitive, so its connected
components have equal size, contrary to
\(2k_1\ne2k_2\). Thus \(r=1\), and the
cross edges form a perfect matching.
(If \(r=0\), \(H=2K_m\) is homogeneous,
the already-listed trivial alternative.)
This avoids the unsupported invocation of Claim 3.9
for the entire \(m\times m\) cross relation.

\section*{Replacement for Claim 3.16}

Suppose we are in Case 2, so that the preceding
double-cover argument has identified \(H\) as the
Taylor double of
\[
 G=\dot\bigcup_{i=1}^s K_{k_i},\qquad
 n=\sum_i k_i,
\]
possibly with the matching between fibre mates.
Let \(T_i\) be the number of triangles of \(H\) through
a vertex over the \(i\)-th clique. Separating
neighbours in its own layer and in the opposite layer
gives
\[
 T_i=\binom{k_i-1}{2}
       +\sum_{j\ne i}\binom{k_j}{2}
       +(k_i-1)(n-k_i)
     =\sum_j\binom{k_j}{2}
       +(k_i-1)(n-k_i-1).
\]
Fibre mates have no common neighbour, so adding the
extra matching creates no triangles. Consequently
\[
 T_i-T_j=(k_i-k_j)(n-k_i-k_j).
\]
If \(s\ge3\), then \(n-k_i-k_j>0\) for every \(i\ne j\).
The \(k_i\)'s are not all equal, so the triangle counts
are not constant. This contradicts vertex-transitivity
of \(H\). Therefore \(s=2\).

\paragraph{Proof boundaries.}
The new arguments above are paper-level deductions
checked independently for logical consistency.
The triangle formula was tested exhaustively for
all 2,024 composition/double cases of base size
\(3\le n\le10\); see
\texttt{checks/section3\_triangle\_double.py}.
They do \emph{not} by themselves prove the earlier
assertion in Case 2 that the graph is indeed a
Taylor double (possibly with a matching), nor the
unconditional version of Lemma 3.12 in the manuscript.
No original manuscript text has been changed.
\end{document}
